\documentclass[conference]{IEEEtran}
\IEEEoverridecommandlockouts
\usepackage{cite}
\usepackage{amsmath,amssymb,amsfonts}
\usepackage{algorithmic}
\usepackage{graphicx}
\usepackage{textcomp}
\usepackage{xcolor}
\def\BibTeX{{\rm B\kern-.05em{\sc i\kern-.025em b}\kern-.08em
    T\kern-.1667em\lower.7ex\hbox{E}\kern-.125emX}}

\begin{document}

\title{Evolving Navigation Policies for FrozenLake-v1 Using Genetic Algorithms:\\
An Empirical Study of Operator Selection, Parameter Sensitivity, and Fitness Shaping}

\author{
\IEEEauthorblockN{[Author 1]}
\IEEEauthorblockA{\textit{[Department]} \\
\textit{[Institution]}\\
[City, Country] \\
[email]}
\and
\IEEEauthorblockN{[Author 2]}
\IEEEauthorblockA{\textit{[Department]} \\
\textit{[Institution]}\\
[City, Country] \\
[email]}
\and
\IEEEauthorblockN{[Author 3]}
\IEEEauthorblockA{\textit{[Department]} \\
\textit{[Institution]}\\
[City, Country] \\
[email]}
}

\maketitle

% ---------------------------------------------------------------
\begin{abstract}
[Placeholder: complete after Methodology and Experimental sections are written.
Approximately 150--250 words summarising motivation, approach, and key findings.]
\end{abstract}

\begin{IEEEkeywords}
genetic algorithm, reinforcement learning, FrozenLake, policy search,
fitness shaping, operator selection, fitness landscape
\end{IEEEkeywords}

% ---------------------------------------------------------------
\section{Introduction}
[Placeholder: present background, problem statement, and objectives.]

% ---------------------------------------------------------------
\section{Literature Review}

Genetic algorithms (GAs) have emerged as a principled approach
to reinforcement learning (RL) policy search, treating the space
of possible mappings from states to actions as a fitness landscape
to be explored through selection, crossover, and mutation~\cite{b1}.
In his survey of evolutionary computation for RL,
Whiteson~\cite{b1} characterises \textit{direct policy search} ---
in which each individual in the evolving population encodes a
complete state-to-action mapping --- as particularly well-suited
to environments with sparse rewards and discrete action spaces.
Tournament selection is identified as providing tunable selection
pressure through a single tournament size parameter~\cite{b1},
and episode-based fitness evaluation is shown to handle sparse
reward signals naturally because the fitness score accumulates
over a complete episode rather than relying on per-step credit
assignment.
This theoretical foundation directly motivates the present work:
our genotype encodes one discrete action per state for the
deterministic 8$\times$8 FrozenLake-v1 gridworld (64 integers,
actions \{Left, Down, Right, Up\}), and we apply tournament
selection with configurable tournament size, enabling an
empirical study of selection pressure.

The FrozenLake environment has been used as a direct benchmark
for GA-based neuroevolution by Faycal and Zito~\cite{b2}.
Their baseline GA evolves the weights of a feedforward network
(16-10-10-4, 300 parameters) to navigate FrozenLake-v1 using
uniform crossover, Gaussian mutation, and elitist selection.
Two operator modifications are introduced: \textit{Multi-step
Mutation} (MSM), which applies several additional
crossover-mutation cycles to elite individuals and adopts
a diversity-based internal fitness measure when the mutated
population underperforms, and \textit{Directed Crossover},
which removes the lowest-magnitude fraction ($\zeta = 0.3$)
of network weights before mating so that only the most
influential weights are inherited~\cite{b2}.
Both modifications achieve convergence in under 150 generations
versus more than 500 for the baseline (Mann-Whitney U,
$p < 0.01$), demonstrating that operator design is a critical
determinant of convergence speed even on a compact gridworld
task.
Although their policy is encoded as neural-network weights
whereas ours uses a tabular direct encoding (one integer action
per state), the core finding --- that crossover and mutation
operator choice substantially alters search efficiency ---
directly motivates our Experiment~1, in which we exhaustively
test all combinations of three crossover methods (one-point,
two-point, uniform) and five mutation methods (random resetting,
swap, insert, scramble, inversion) across two replacement
strategies ($\mu{+}\lambda$ and $\mu$,$\lambda$), validated
with a pairwise $t$-test between the top two configurations.

The interaction between operator choice, operator rates, and
population parameters has been studied systematically in the
context of hybrid RL-assisted evolutionary algorithms.
Song et al.~\cite{b3} provide a comprehensive taxonomy of
reinforcement learning-assisted evolutionary algorithms (RL-EA),
classifying integration strategies into five categories:
solution generation, learnable objective functions,
algorithm/operator/sub-population selection, parameter
adaptation, and other strategies.
A key finding from this survey is that adaptive operator
selection consistently reduces premature convergence and
improves final solution quality relative to fixed-operator
configurations~\cite{b3}.
The survey further establishes that representing population
state as a continuous feature vector --- including mean fitness,
best fitness, fitness standard deviation, and iteration
progress --- enables an RL agent to select optimal operator
combinations for the current stage of evolution.
This provides the theoretical motivation for our Experiment~2,
which systematically quantifies the sensitivity of fitness to
mutation rate, crossover rate, population size, and tournament
size through grid-search heatmaps, generating the empirical
data that an RL-EA operator-selection agent would use as a
learning signal.

Adaptive operator and parameter selection in GA has been
demonstrated concretely by Ma et al.~\cite{b4}, who propose
DRL-A-GA for the flexible job shop scheduling problem (FJSP).
A double DQN agent observes a four-dimensional continuous
population state vector (normalised mean fitness, best fitness,
fitness standard deviation, and iteration progress) and
selects both the operator type --- from four crossover
operators and four domain-specific mutation operators ---
and the rates (crossover rate $\in \{0.8, 0.9, 1.0\}$,
mutation rate $\in \{0.1, 0.2, 0.3\}$) at each
generation~\cite{b4}.
An ablation study confirms that adapting both operator type
and rates simultaneously is necessary: either alone yields
significantly worse results, with the full DRL-A-GA achieving
the lowest relative percentage deviation (16.73\%) across
10 Brandimarte FJSP benchmark instances.
The parallel to our project is direct: our Experiment~1
evaluates all 30 operator-combination configurations
exhaustively and applies statistical testing to identify
the best combination, while our Experiments~2a and~2b
use heatmaps over mutation rate vs.\ crossover rate and
population size vs.\ tournament size, respectively ---
precisely the parameter axes that DRL-A-GA learns to
navigate adaptively~\cite{b4}.

A critical design choice in any GA applied to sparse-reward
navigation is the fitness function.
Zhu et al.~\cite{b5} address this challenge in the context
of partially observable gridworld RL by proposing potential-based
reward automata optimised via genetic local search (GLS-PRA).
Their framework augments the sparse terminal reward with a
potential-based shaping term $\gamma\phi(e') - \phi(e)$,
where $\phi$ is a value function computed over reward automaton
states via value iteration, and prove (Theorem~4.4) that this
shaping is policy-preserving regardless of the potential
function chosen~\cite{b5}.
Empirically, potential-based shaping increases cumulative
reward by up to 39.6\% in a key navigation gridworld domain
(from 95.40 to 133.22) while substantially reducing training
variance.
Our fitness function employs an analogous mechanism: the
Manhattan-distance component
$w_m / (1 + d_\mathrm{Manhattan}(\text{final state}, \text{goal}))$
with $w_m = 10.0$ acts as a surrogate potential function,
assigning non-zero fitness to every episode regardless of
whether the agent reaches the goal, and thereby guiding
evolution on the sparse-reward FrozenLake landscape.
A step penalty ($w_s = 0.1$ per step) further discourages
inefficient trajectories, and a hole penalty ($w_h = 50.0$)
mirrors the obstacle-avoidance shaping used in the gridworld
benchmarks of Zhu et al.~\cite{b5}.
Our Experiment~2c directly tests the sensitivity of the GA
to these fitness weight choices by comparing the standard
configuration against a heavy-hole-penalty variant
($w_h = 200.0$) and a heavy-step-penalty variant
($w_s = 2.0$).

The theoretical underpinning of why mutation rate, tournament
size, and population size interact non-trivially in GA search
is provided by fitness landscape theory.
Srivastava et al.~\cite{b6} review fitness landscapes as
mappings from genotypes to fitness and formalise
\textit{navigability} as the probability of reaching the
global fitness peak via random mutation and selection from
an arbitrary starting genotype.
Their analysis shows that rugged landscapes --- characterised
by many local maxima and sign epistasis, where the fitness
effect of a mutation reverses depending on the genetic
background --- require a careful balance between exploration
(maintained by moderate mutation and low selection pressure)
and exploitation (driven by strong selection and low
mutation) to navigate successfully~\cite{b6}.
In our context, the FrozenLake-v1 policy space is a discrete
fitness landscape of $4^{64} \approx 3.4 \times 10^{38}$
possible genotypes (one of four actions per each of 64 states).
Its ruggedness is apparent: policies that differ by a single
gene value at a critical junction state --- for instance,
the state immediately adjacent to a hole --- can differ
drastically in fitness, an instance of sign epistasis as
characterised by Srivastava et al.~\cite{b6}.
Their finding that population diversity, maintained via
lower selection pressure or higher mutation rate, aids
navigation of rugged landscapes directly informs the
hypotheses for our Experiments~2a and~2b: we expect
intermediate mutation rates and moderate tournament sizes
to yield the best trade-off between exploration and
convergence, and we expect this optimum to shift when
the fitness landscape is made more rugged by increasing
grid size (Experiment~3, scalability study across
8$\times$8, 16$\times$16, and 32$\times$32 maps).

Taken together, these six studies establish the key
dimensions of our investigation.
Whiteson~\cite{b1} and Faycal and Zito~\cite{b2} confirm
that GA-based direct policy search is both theoretically
justified and empirically effective on gridworld RL tasks,
with operator choice being a primary driver of convergence
speed.
Song et al.~\cite{b3} and Ma et al.~\cite{b4} demonstrate
that systematic analysis of operator and parameter
combinations is foundational to understanding EA performance
and provides the data required to build adaptive systems.
Zhu et al.~\cite{b5} provide theoretical and empirical
support for the potential-based shaping component in our
multi-component fitness function.
Srivastava et al.~\cite{b6} supply the landscape-theoretic
framework for interpreting the observed sensitivities to
mutation rate, population size, and selection pressure
across our parameter-sensitivity and scalability experiments.
Our work contributes an empirical study that spans all three
dimensions simultaneously --- operator selection, parameter
sensitivity, and fitness shaping --- on a canonical
deterministic gridworld task, using a tabular direct-encoding
GA with statistical validation and scalability analysis.

% ---------------------------------------------------------------
\section{Methodology}
[Placeholder: describe GA representation, fitness function,
operators, and experimental setup.]

% ---------------------------------------------------------------
\section{Experimental Work}
[Placeholder: present and discuss numerical results from
Experiments 1, 2a, 2b, 2c, and 3. Compare with literature.]

% ---------------------------------------------------------------
\section{Conclusion and Discussion}
[Placeholder: summarise findings, reflect on implications,
and suggest future work.]

% ---------------------------------------------------------------
\section*{AI Usage Statement}
[Placeholder: detail which generative AI tools were used,
the scope of their application, and how outputs were verified.]

% ---------------------------------------------------------------
\section*{Data Availability}
The source code, configuration file, and experimental results
are available at: [repository link to be added].

% ---------------------------------------------------------------
\begin{thebibliography}{00}

\bibitem{b1}
S. Whiteson, ``Evolutionary computation for reinforcement
learning,'' in \textit{Reinforcement Learning: State-of-the-Art},
Springer, 2012, pp.~325--355.

\bibitem{b2}
T. Faycal and C. Zito, ``Direct mutation and crossover in
genetic algorithms applied to reinforcement learning tasks,''
\textit{arXiv preprint} arXiv:2201.04815, 2022.

\bibitem{b3}
Y. Song, Y. Wu, Y. Guo, R. Yan, P.~N. Suganthan, Y. Zhang,
W. Pedrycz, S. Das, R. Mallipeddi, O.~S. Ajani, and Q. Feng,
``Reinforcement learning-assisted evolutionary algorithm:
A survey and research opportunities,''
\textit{Swarm Evol. Comput.}, vol.~86, p.~101517, 2024.

\bibitem{b4}
J. Ma, W. Gao, and W. Tong, ``A deep reinforcement learning
assisted adaptive genetic algorithm for flexible job shop
scheduling,'' \textit{Eng. Appl. Artif. Intell.}, vol.~149,
p.~110447, 2025.

\bibitem{b5}
Z. Zhu, Z. Chen, C. Zhu, W. Si, and F. Wang, ``Optimizing
potential-based reward automata in partially observable
reinforcement learning using genetic local search,''
\textit{Eng. Appl. Artif. Intell.}, vol.~169, p.~114054, 2026.

\bibitem{b6}
M. Srivastava, C. Bank, J. Krug, and S.~G. Das, ``Evolution
as fitness landscape navigation: Concepts, measures, and
emerging questions,'' \textit{arXiv preprint}
arXiv:2604.17036, 2026.

\end{thebibliography}

\end{document}
